\documentclass[12pt]{article}
\usepackage[utf8]{inputenc}
\usepackage[T1]{fontenc}
\usepackage{amsmath, amssymb, geometry, booktabs}
\usepackage{enumitem}
\geometry{margin=1in}
\setlength{\parindent}{0pt}
\renewcommand{\arraystretch}{1.2}

\begin{document}



\section*{Problem 1 (13 marks) }

Two countries, \textbf{Home (H)} and \textbf{Foreign (F)}, produce two goods: \textbf{textiles (T)} (labour-intensive) and \textbf{steel (S)} (capital-intensive).  
Both countries have the same technologies and homothetic preferences.

\begin{center}
\begin{tabular}{lccc}
\toprule
\textbf{Country} & \textbf{Capital $K$} & \textbf{Labour $L$} & \textbf{$K/L$ ratio} \\
\midrule
Home & 300 & 900 & 0.33 \\
Foreign & 800 & 800 & 1.00 \\
\bottomrule
\end{tabular}
\end{center}

Assume perfect competition, identical production functions, and full employment.

\begin{enumerate}[label=(\arabic*)]
\item Identify which country is relatively labour-abundant and which is capital-abundant.  
\item According to the Heckscher--Ohlin theorem, which good will each country export under free trade?  
\item Using the Stolper--Samuelson theorem, explain how opening to trade affects the real wage and the real rental rate in each country.  
\item Suppose the world relative price of textiles $(P_T/P_S)$ rises by 10\%.  
Explain how this change affects factor incomes (wages and rents) and the output composition in the Home country.  
\item If trade continues over time, what mechanism could lead to \textbf{factor price equalization} between the two countries?  
\end{enumerate}





\end{document}
